% Copyright (c) 2026 Sergey Zhunkov.
% SPDX-License-Identifier: CC-BY-4.0
% https://creativecommons.org/licenses/by/4.0/
%
% Shared by the seven-note and pentatonic CAGED maps.
\usepackage[margin=0pt]{geometry}
\usepackage{fontspec}
\usepackage{babel}
\usepackage{etoolbox}
\usepackage{tikz}
\usetikzlibrary{babel}
\usepackage{luacode}
\usepackage{hyperref}

\providecommand{\ChartLanguage}{en}
\newif\ifChartRussian
\ifdefstring{\ChartLanguage}{ru}{%
  \ChartRussiantrue
  \babelprovide[import,main]{russian}
  \babelprovide[import]{english}
}{%
  \ifdefstring{\ChartLanguage}{en}{%
    \babelprovide[import,main]{english}
    \babelprovide[import]{russian}
  }{%
    \PackageError{guitar-scale-maps}{Unsupported ChartLanguage '\ChartLanguage'}
      {Set ChartLanguage to en or ru.}%
  }%
}

\setmainfont{DejaVuSansCondensed.ttf}[
  BoldFont=DejaVuSansCondensed-Bold.ttf,
  Ligatures=NoCommon,
  RawFeature={-kern,-clig}
]
\pagestyle{empty}

\definecolor{ChartInk}{RGB}{23,43,50}
\definecolor{ChartMuted}{RGB}{70,91,99}
\providecommand{\ChartAccentHex}{195755}
\definecolor{ChartAccent}{HTML}{\ChartAccentHex}
\definecolor{ChartRule}{RGB}{197,208,209}
\definecolor{ChartStripe}{RGB}{244,247,247}
\colorlet{ChartHeader}{ChartAccent!13!white}
\colorlet{ChartHeading}{ChartAccent!33!ChartInk}

\newcommand{\Localize}[2]{\ifChartRussian#1\else#2\fi}

\begin{luacode*}
CagedMap = {}
-- Sounding MIDI pitches, high to low (strings 1-6), as on the fretboard map.
CagedMap.open_strings = {
  {64, "E4"}, {59, "B3"}, {55, "G3"}, {50, "D3"}, {45, "A2"}, {40, "E2"},
}

-- CAGED letters come from transposing the corresponding open major chord
-- forms to C. Each window includes all scale notes on its inclusive frets.
-- A and E include the fret below their chord form. D spans five frets to
-- include the fret below its form and the 2nd-string tonic at fret 13.
-- The C form is shown at the octave so its full window has frets above 0.
-- A natural minor uses the same windows and labels degrees from A instead.
local e_position = {"E", 7, 10}
CagedMap.positions = {{"A", 2, 5}, {"G", 5, 8}, e_position, {"D", 9, 13}, {"C", 12, 15}}
local positions = CagedMap.positions
local overview_frets = 24
local example_shift = 2

function CagedMap.by_pitch(tonic, intervals, values)
  local by_pitch = {}
  for index, interval in ipairs(intervals) do
    by_pitch[(tonic + interval) % 12] = values[index]
  end
  return by_pitch
end

function CagedMap.print_grid_notes(degrees_by_pitch, first_fret, last_fret)
  for string_number, string_data in ipairs(CagedMap.open_strings) do
    local open_pitch, open_note = table.unpack(string_data)
    tex.sprint(string.format([[\StringLabel{%d}{%s}]], string_number, open_note))
    for fret = first_fret, last_fret do
      local degree = degrees_by_pitch[(open_pitch + fret) % 12]
      if degree then
        tex.sprint(string.format([[\FretDegree{%d}{%d}{%s}]],
          string_number, fret, degree))
      end
    end
  end
end

function PrintPositions(key, example)
  for index, position in ipairs(positions) do
    local name, first, last = table.unpack(position)
    tex.sprint(string.format([[\PositionCard{%d}{%s}{%d}{%d}{%s}]],
      index, name, first, last, key))
  end
  local name, first, last = table.unpack(e_position)
  tex.sprint(string.format([[\ExampleCard{%s}{%d}{%d}{%s}]],
    name, first + example_shift, last + example_shift, example))
end

function PrintCoverage()
  -- Include the tail of the previous cycle and the start of the next.
  -- Consecutive windows alternate between two lanes so their overlap is visible.
  for octave = -1, 1 do
    for index, position in ipairs(positions) do
      local name, first, last = table.unpack(position)
      first = math.max(1, first + 12 * octave)
      last = math.min(overview_frets, last + 12 * octave)
      if first <= last then
        local lane = (octave * #positions + index) % 2
        tex.sprint(string.format([[\CoverageBand{%s}{%d}{%d}{%d}]],
          name, first, last, lane))
      end
    end
  end
end
\end{luacode*}

% Coordinates and font sizes are in PDF points, measured from the bottom left.
\newcommand{\ChartText}[6][base west]{%
  \node[anchor=#1,inner sep=0pt,outer sep=0pt,text=#5,
    font=\fontsize{#4bp}{#4bp}\selectfont] at ({#2},{#3}) {#6};%
}
\newcommand{\ChartLeft}{34}
\newcommand{\ChartRight}{808}
\newcommand{\OverviewGridLeft}{94}

% Numbered fret cells increase to the right, with string 1 at the top.
\newcommand{\FretEdge}[1]{\GridLeft+\FretWidth*(#1-\FirstFret)}
\newcommand{\FretCenter}[1]{\FretEdge{#1+.5}}
\newcommand{\StringEdge}[1]{\GridTop-19-17*(#1)}
\newcommand{\StringCenter}[1]{\StringEdge{#1-.5}}

\newcommand{\DegreeMarker}[3]{%
  \ifstrequal{#3}{1}{%
    \fill[ChartHeading,rounded corners=1bp]
      ({#1-7},{#2-6.7}) rectangle ({#1+7},{#2+6.7});%
    \ChartText[center]{#1}{#2}{10}{white}{\bfseries 1}%
  }{%
    \draw[ChartMuted,line width=.55bp,fill=white]
      ({#1},{#2}) circle[radius=6.7bp];%
    \ChartText[center]{#1}{#2}{9.5}{ChartInk}{#3}%
  }%
}
\newcommand{\FretDegree}[3]{%
  \DegreeMarker{\FretCenter{#2}}{\StringCenter{#1}}{#3}%
}

\newcommand{\StringLabel}[2]{%
  \ChartText[base east]{\GridLeft-7}{\StringCenter{#1}-3}{9}{ChartMuted}%
    {\bfseries#1\enspace\mdseries#2}%
}

% Grid arguments are left edge, top edge, width, first/last fret, and scale key.
\newcommand{\DrawGrid}[6]{%
  \begin{scope}
    \def\GridLeft{#1}\def\GridTop{#2}\def\FirstFret{#4}%
    \pgfmathsetmacro{\GridRight}{#1+#3}%
    \pgfmathsetmacro{\FretWidth}{(#3)/(#5-#4+1)}%
    \fill[ChartHeader] (\GridLeft,\GridTop) rectangle (\GridRight,{\StringEdge{0}});
    \foreach \row in {2,4,6} {%
      \fill[ChartStripe] (\GridLeft,{\StringEdge{\row-1}})
        rectangle (\GridRight,{\StringEdge{\row}});
    }%
    \foreach \fret in {#4,...,#5} {%
      \draw[ChartRule,line width=.50bp] ({\FretEdge{\fret}},\GridTop)
        -- ({\FretEdge{\fret}},{\StringEdge{6}});
      \ChartText[base]{\FretCenter{\fret}}{\GridTop-12.5}{10}{ChartHeading}{\bfseries\fret}%
    }%
    \foreach \row in {0,...,5} {%
      \draw[ChartRule,line width=.50bp] (\GridLeft,{\StringEdge{\row}})
        -- (\GridRight,{\StringEdge{\row}});
    }%
    \draw[ChartRule,line width=.60bp] (\GridLeft,\GridTop)
      rectangle (\GridRight,{\StringEdge{6}});
    \directlua{PrintGridNotes("#6", #4, #5)}%
  \end{scope}%
}

\newcommand{\DrawHeading}[2]{%
  \ChartText{\ChartLeft}{542}{28}{ChartInk}{\bfseries #1}%
  \ChartText[base east]{\ChartRight}{542}{12}{ChartHeading}{\bfseries #2}%
  \draw[ChartRule,line width=.65bp] (\ChartLeft,529) -- (\ChartRight,529);
}

\newcommand{\CoverageBand}[4]{%
  \pgfmathsetmacro{\BandLeft}{\OverviewGridLeft+(\ChartRight-\OverviewGridLeft)/24*(#2-1)+.8}%
  \pgfmathsetmacro{\BandRight}{\OverviewGridLeft+(\ChartRight-\OverviewGridLeft)/24*#3-.8}%
  \pgfmathsetmacro{\BandBottom}{422+17*#4}%
  \draw[draw=ChartHeading,fill=ChartHeader,line width=.55bp,rounded corners=1.5bp]
    (\BandLeft,\BandBottom) rectangle (\BandRight,{\BandBottom+13.5});
  \ChartText[center]{(\BandLeft+\BandRight)/2}{\BandBottom+6.75}{9.5}{ChartHeading}{\bfseries #1}%
}

\newcommand{\ScaleKeyCell}[3]{%
  \pgfmathsetmacro{\KeyX}{\KeyStart+35*(#1-1)}%
  \ChartText[base]{\KeyX}{\KeyTop}{\KeySize}{ChartHeading}{\bfseries #2}%
  \ChartText[base]{\KeyX}{\KeyTop-14}{\KeySize}{ChartInk}{#3}%
}
\newcommand{\ScaleKey}[4]{%
  \begin{scope}
    \def\KeyStart{#1}\def\KeyTop{#2}\def\KeySize{#3}%
    \ChartText[base east]{#1-23}{#2}{#3}{ChartMuted}{\Localize{Ступень}{Degree}}%
    \ChartText[base east]{#1-23}{#2-14}{#3}{ChartMuted}{\Localize{Нота}{Note}}%
    \directlua{PrintScaleKey("#4")}%
  \end{scope}%
}

\newcommand{\DrawCard}[6]{%
  \begin{scope}[shift={(#1,#2)}]
    \fill[ChartHeader] (0,-24) rectangle (248,0);
    \draw[ChartRule,line width=.60bp] (0,0) rectangle (248,-161);
    \ChartText{12}{-16}{12}{ChartHeading}{\bfseries #3}%
    \ChartText[base east]{236}{-16}{10.5}{ChartInk}%
      {\Localize{Лады}{Frets} #4–#5}%
    \DrawGrid{51}{-32}{185}{#4}{#5}{#6}%
  \end{scope}%
}
\newcommand{\PositionCard}[5]{%
  \pgfmathsetmacro{\CardX}{\ChartLeft+263*mod(#1-1,3)}%
  \pgfmathsetmacro{\CardY}{452-173*floor((#1-1)/3)}%
  \DrawCard{\CardX}{\CardY}{\Localize{Форма}{Form} #2}{#3}{#4}{#5}%
}
\newcommand{\ExampleCard}[4]{%
  \DrawCard{560}{279}{\ExampleTitle\space· #1}{#2}{#3}{#4}%
}

\newcommand{\ReadmeReference}[1]{%
  https://github.com/zhunkoff/music-reference/blob/main/\Localize{README.ru.md}{README.md}\#ref-#1%
}
\newcommand{\DrawFooter}{%
  \draw[ChartRule,line width=.65bp] (\ChartLeft,50) -- (\ChartRight,50);
  \ChartText[base east]{\ChartRight}{19}{9}{ChartMuted}{%
    © 2026 \href{https://github.com/zhunkoff/music-reference}{Sergey Zhunkov} ·
    \href{https://creativecommons.org/licenses/by/4.0/}{CC BY 4.0}}%
}

\newcommand{\PrintChartPage}[2]{%
  \pdfbookmark[0]{#1}{scale-maps-\thepage}%
  \AddToHookNext{shipout/background}{%
    \put(0,-\paperheight){\begin{tikzpicture}[overlay,x=1bp,y=1bp]#2\end{tikzpicture}}%
  }%
  \null\newpage
}
